\subsection{The actual geometric extension selects the integral phase}

\begin{statement}{B.40}{The exact local action class and the integral Ross refinement}{Proved}
Keep the tangent at $C_0$, the primitive vector $w$, and all
coefficient conventions of B.24 and B.35--B.38.
In the thirteen Smith coordinates of B.36 the actual class is
\[
 [\eta]_{\mathbb Q_2}
  =(0,0,1,2,0;\ 2,5,9,6,9,24;\ 57,36)
  =b+6h_9.                                                     \tag{B.40.1}
\]
The coordinate moduli are
\[
 (3,3,3,3,3;\ 9,9,27,27,27,27;\ 81,81).
\]
In particular the two intrinsic finite detectors of B.38--B.39
have the exact values
\[
          \lambda_{81}([\eta])=57\pmod {81},\qquad
          \lambda_{27}([\eta])=19\pmod {27}.                    \tag{B.40.2}
\]
Here $\lambda_{27}=\lambda_{81}|_{\mathcal J}/3$ uses the
actual $\mathcal J$-valued commutator lift, as in B.38--B.39.

The actual geometric action supplies an additional integral
condition that distinguishes the three formal refinements in
B.38. Exactly one satisfies that condition. Together with
B.20 and B.37, this determines the requested rational
coefficient and integral refinement: the rational coefficient
is $-1$ relative to the specified Ross cup edge, and the
remaining integral phase is the one selected by (B.40.1).
The Ross symbol's real-regulator coordinate $-3/2$ in B.19
is a different normalization and is not replaced by $-1$.
No numerical value of a three-adic regulator is asserted.

Thus the proposed inherited quotient of order $3^{30}$ in
task B.4 cannot perform the requested separation, by B.22
and B.38. The intrinsic quotient of B.39 does perform it,
and its invariant is evaluated by (B.40.2). This resolves
the integral ambiguity without assigning an unproved
global Frobenius word or treating a finite three-adic
truncation as an exact three-adic number.
\end{statement}
\begin{proof}
\par\smallskip\noindent\textbf{The geometric condition missing from a commutation identity.}\enspace 
Write $M=A(-1)$ and $j=\operatorname{ad}_w:V\hookrightarrow M$.
Its image is primitive by B.35, so
\[
                         Q=M/jV
\]
is a free lattice of rank $320$.
Let $L_w$ be the inverse image of $\mathbb Z_3w$ in
$(\gamma_2N/\gamma_4N)(-1)$.
The extension
\[
 0\longrightarrow Q\longrightarrow \overline L_w
       :=L_w/jV\longrightarrow\mathbb Z_3\longrightarrow0
                                                               \tag{B.40.3}
\]
has a canonical action of the geometric automorphism group
$\mathcal A=E[3]\rtimes\mu_6$, of order $54$.
Indeed inner automorphisms by $N$ act trivially on $M$ and
send a lift of $w$ to that lift plus $[v,w]$ for some $v\in V$.
They therefore act trivially on (B.40.3).
Every element of $\mathcal A$ has compact determinant one.
Consequently the quotient $\mathbb Z_3$ in (B.40.3) has
trivial geometric action, and choices of paths used to
base the spatial automorphisms no longer affect this module.

Over $K=\mathbb Q(\zeta_3)$ these geometric automorphisms
are defined over the arithmetic base. Their actions on
(B.40.3) commute with the arithmetic action. Denote the
resulting geometric extension class by
\[
                   \xi\in H^1(\mathcal A,Q).
\]
This is actual geometric data, not just the assertion that
an arithmetic class is invariant under $\mathcal A$.

At two pass to the unramified quadratic extension with
residue field $\mathbb F_4$. Let $F$ denote its arithmetic
Frobenius on $Q$ and put $D=F-I$.
The rational eigenvalues of $D$ are among
$-3,3,-9,15$, so $D$ is invertible over $\mathbb Q_3$.
Choose a lift $\omega$ of $w$ and write
\[
       a\omega=\omega+t_a,\qquad
       F\omega=\omega+e
\]
in the twisted extension, where $e$ represents the image
of $[\eta](F_4)$ in $Q/DQ$.
Commutativity in (B.40.3) gives the literal identity
\[
                   (a-I)e=Dt_a.
\]
It follows that
\[
 \xi=\bigl[a\longmapsto(a-I)D^{-1}e\bigr]
                  \quad\hbox{in }H^1(\mathcal A,Q).           \tag{B.40.4}
\]
Changing $e$ by $Dx$ changes the right side by the
coboundary of $x\in Q$, as required.

For a torsion difference represented by
$y\in Q_{\mathbb Q_3}/Q$, the corresponding change of
geometric extension is $\partial_{\mathcal A}(y)$.
It vanishes exactly when
\[
                         y\in Q+Q_{\mathbb Q_3}^{\mathcal A}.
                                                               \tag{B.40.5}
\]
In particular, for a residual class $z/3$ from B.38,
this condition says that $\bar z\in Q/3Q$ is the reduction
of an element of the integral invariant lattice $Q^{\mathcal A}$.
This follows directly by multiplying (B.40.5) by three;
no rational averaging is substituted for integral invariance.

The invariant lattice has rank four. There is a small
exact certificate. In the adapted quotient Hall basis
of B.37 let $S=\sum_{a\in\mathcal A}a$.
Its rows $46,51,80,81$ and columns $4,13,40,66$, numbered
from zero, give the minor
\[
 \begin{pmatrix}
 -1&2&0&0\\ -6&6&0&1\\ -2&2&0&0\\ -1&-2&1&1
 \end{pmatrix},\qquad \det=2.                                 \tag{B.40.6}
\]
The exact identities $S^2=54S$ and
$\operatorname{tr}S=216$ show that its rational image
has rank four. The unit minor proves that these four
columns span a primitive lattice, hence the entire
$Q^{\mathcal A}$ over $\mathbb Z_3$.
The row with coefficients $(2,2,2,1)$ in quotient
positions $(51,80,81,82)$ kills this lattice modulo
three and takes value one on the oriented residual
numerator of B.38. Thus the nonzero residual classes
change $\xi$.
Both the complete integer matrix and this four-entry
certificate are saved by
\texttt{c\_\allowbreak geometric\_\allowbreak extension\_\allowbreak test.py}; an independent
Smith calculation is
\texttt{b\_\allowbreak geometric\_\allowbreak invariants.py}.

\par\smallskip\noindent\textbf{The actual geometric extension from a different parent cover.}\enspace 
Use the multiplication-three cover
\[
                [3]:E-E[3]\longrightarrow E-\{O\}.           \tag{B.40.7}
\]
The parent in (B.40.7) is a once-punctured elliptic curve;
it is not the tripod parent of B.22 and B.32.
Over $\mathbb C$ its topological group is free on two
generators. The index-nine kernel associated with (B.40.7)
is
\[
 F_{\rm top}=\langle a,b\rangle,\qquad
 N_{\rm top}=\ker\bigl(F_{\rm top}\longrightarrow(\mathbb Z/3)^2\bigr).
                                                               \tag{B.40.8}
\]
These are discrete topological groups. Riemann existence
identifies their profinite completions with the geometric
\'{e}tale groups of the displayed cover. Taking pro-three
completions identifies $N_{\rm top}^{(3)}$ with the earlier
geometric group $N$ and with the kernel in
$F_{\rm top}^{(3)}$. The index is a power of three:
for any normal subgroup of three-power index in
$N_{\rm top}$, its core under the nine parent cosets
still has three-power index in $F_{\rm top}$.
Thus the induced pro-three topology on $N_{\rm top}$
is its full pro-three topology. All following word
computations take place in the discrete groups and
extend continuously to these completions.

In this paragraph words concatenate chronological paths
from left to right. Choose the positive peripheral word
$\delta=aba^{-1}b^{-1}$ by cutting the torus along the
two oriented generators. The nine cusp loops are
$t_c\delta t_c^{-1}$, where
\[
 t_{i,j}=a^ib^j,\qquad 0\le i,j\le2.
\]
The deck labels of the eight nonzero cusps in the order
of B.24 are
\[
 (1,0),(2,2),(0,1),(2,0),(1,1),(0,2),(2,1),(1,2).
                                                               \tag{B.40.9}
\]
The Schreier generator for an edge $(c,x)$ is
$t_cx\,t_{c+x}^{-1}$, with the trivial tree edges omitted.
The ten resulting words, in the executed order, are
\[
\begin{aligned}
 &bab^{-1}a^{-1},\quad b^2ab^{-2}a^{-1},\quad b^3,\\
 &abab^{-1}a^{-2},\quad ab^2ab^{-2}a^{-2},\quad ab^3a^{-1},\\
 &a^3,\quad a^2bab^{-1},\quad a^2b^2ab^{-2},\quad a^2b^3a^{-2}.
\end{aligned}                                                  \tag{B.40.10}
\]
Rewriting each of these words gives its own single free
generator. This verifies the marking as well as the
rank-ten count.

Let $e_c$ be the abelianized peripheral word at $c$.
The two transfer vectors, in the basis (B.40.10), are
\[
\begin{aligned}
 u_F&=(1,1,0,1,1,0,1,1,1,0)^t,\\
 v_F&=(0,0,1,0,0,1,0,0,0,1)^t.
\end{aligned}                                                  \tag{B.40.11}
\]
With the eight-column matrix $E=(e_c)$ and the literal
matrix $U$ of B.24, the change of basis
\[
       P=\left(\frac{(u_F,v_F)-EU}{3},\ E\right)              \tag{B.40.12}
\]
has integer entries and determinant $-1$.
It sends the transfer vectors to the specified
$u=(3,0,U_1)^t$ and $v=(0,3,U_2)^t$.
Thus it identifies the marked abelianization with the
actual integral lattice of B.24, including all boundary
vectors. Conjugation by $a$ and $b$ gives exactly its
two translation matrices, not merely their reductions
modulo three.

For the two elliptic automorphisms fixing $O$ use
\[
\begin{array}{c|cc}
 &a&b\\ \hline
 \sigma&aba^{-1}&b^{-1}a^{-1}\\
 \iota&aba^{-1}b^{-1}a^{-1}&ab^{-1}a^{-1}.
\end{array}                                                     \tag{B.40.13}
\]
Both preserve $\delta$ literally. Their compact matrices
are respectively $R_3$ and $-I$, and their permutations
of the nine cusp classes are the actual permutations
of these Fermat automorphisms.
They satisfy $\sigma^3=\operatorname{Int}(\delta)$ and
$\iota^2=\operatorname{Int}(\delta)$ as based word
automorphisms. They are not incorrectly declared to
have orders three and two before passing to the outer module.

Here the identification with the actual geometric action
uses the parent torus, not the false principle that an
arbitrary automorphism of a rank-ten free group is
determined by its abelianization.
For a once-punctured torus its oriented mapping class
group is $\operatorname{SL}_2(\mathbb Z)$, and
$\operatorname{Out}(F_2)=\operatorname{GL}_2(\mathbb Z)$;
see \citet[\S2.2.4, pp.~54--55; \S8.2.7,
pp.~234--235]{BFarbMargalit2012}.
The actual CM map and the word map in (B.40.13) therefore
have the same outer action on the parent group.
An inner discrepancy which preserves the chosen origin
cusp in the cover has trivial deck translation and lies
in $N$. Its action is killed in (B.40.3).
The same argument applies to $\iota$.
The two deck translations together with these two maps
therefore give the actual $\mathcal A$-action on
(B.40.3).

\par\smallskip\noindent\textbf{The sign from chronological paths to fibre automorphisms.}\enspace 
This comparison must retain the order of path composition.
If one first follows a path $\gamma$ and then a path
$\gamma'$, its transport on a fibre is
$M_{\gamma'}\circ M_\gamma$.
Thus the chronological word model used in
(B.40.8)--(B.40.13) maps anti-multiplicatively to
automorphisms of the fibre functor.
It induces the identity on $V$: on the cover
$z\mapsto z^{3^r}$ of $\mathbb G_m$ the positive circle
sends $1$ to the positive $3^r$-th root of unity, and
on an Abel cover the positive period loop sends zero
to its positive division point.
On the associated graded Lie algebra this same
anti-map has signs $+,-,+$ in degrees one, two, and three.
In particular
\[
       B_{\rm fibre}=-B_{\rm chronological}                  \tag{B.40.14}
\]
for the quadratic part of a generator action.
Equivalently every image word is reversed before its
quadratic fibre coordinates are computed.

There is also a based tangent check of this sign.
The chronological map $\sigma$ in (B.40.13) has cube
$\operatorname{Int}(\delta)$.
After fibre transport this becomes
$\operatorname{Int}(M_\delta^{-1})$.
For a positive one-third tangent rotation, the based
map is $p^{-1}\sigma_*p$; its third power is conjugation
by the inverse of the positive peripheral transport.
Thus the word calculation agrees with the tangent
normalization of B.35--B.36 independently of any
candidate comparison.
In particular the positive Kummer identity in B.36 is
unchanged. If $\alpha=p^{-1}F(p)$, then
$TF=\operatorname{Int}(\alpha)FT$, which gives
$(T-I)\eta=[\alpha,w]+Dt$ with the published signs.

\par\smallskip\noindent\textbf{Integral word computation of the geometric class.}\enspace 
In class-two logarithmic coordinates write a word as
$(x,z)$, where $z\in\Lambda^2V$.
The multiplication rule
\[
 (x,z)(y,t)=(x+y,z+t+\tfrac12x\wedge y)
                                                               \tag{B.40.15}
\]
is integral at three.
The program stores $2z$, and updates it by $x\wedge y$
for each letter. Reversing a word leaves its linear
coordinate unchanged and negates $2z$; this is checked
for all forty generator images of (B.40.13) and the
two deck translations.
No class-three Baker--Campbell--Hausdorff expression
with a denominator divisible by three is used.

Let $C_a$ and $B_a$ be the linear and quadratic fibre
matrices of one of these four actions, transported by
(B.40.12). Choose the pure integral commutator lift
of B.38. Its cubic displacement is
\[
 t_a=[B_au,C_av]+[C_au,B_av].                                  \tag{B.40.16}
\]
Cubic terms in generator images cannot enter this
degree-three expression. The image of (B.40.16) in
$Q$ is computed by the primitive Hall deletion of
B.37. Enumerating the group generated by the four
degree-one matrices gives exactly $54$ elements.
For every repeated group element the computed
displacement agrees, verifying all path relations
in the outer module.

For a finite group cocycle $t$ put
\[
              h=-\frac1{54}\sum_{a\in\mathcal A}t_a.
                                                               \tag{B.40.17}
\]
The cocycle identity gives $t_a=(a-I)h$.
This computes the actual geometric extension as the
rational coset $h\bmod Q$. The saved numerator is
$27h=-\frac12\sum t_a$, so its only denominator is
the explicit power $27$.

\par\smallskip\noindent\textbf{The compact comparison does not hide a new phase.}\enspace 
The Schreier calculation uses the CM matrix $R_3$;
B.36 uses the compatible CM matrix $S$ in the local
quaternion basis.
Choose a strict symplectic intertwiner
\[
 B\equiv I\pmod3,\qquad \det B=1,\qquad SB=BR_3.
                                                               \tag{B.40.18}
\]
One explicit choice begins with $B_0=(e_1,Se_1)$ and
multiplies it by the unique scalar $s\equiv1\pmod3$
such that $s^2\det B_0=1$.
The full lattice comparison is
\[
 K(B)=\begin{pmatrix}
 B&0\\ (UB-U)/3&I_8
 \end{pmatrix}.                                                \tag{B.40.19}
\]
It fixes the positive boundary Tate frame and preserves
$w$. For precision, fix compatible positive roots of
unity in every three-power Kummer cover, extending the
specified $\zeta_3$. These give the boundary frame over
$\mathbb Z_3(1)$, not merely its reduction modulo three.
The compact Weil pairing uses this same Tate frame.
The parent peripheral commutator fixes the corresponding
compact intersection pairing, so a comparison preserving
these marked data has determinant exactly one in
$\mathbb Z_3$. This is the strict symplectic normalization
in (B.40.18), not an inference from determinant one
modulo three alone.
The rational compact and boundary summands have different
Frobenius-four eigenvalues, so this is the unique full
degree-one comparison with the indicated compact and
boundary parts.

An actual comparison can differ from (B.40.18) by a
strict symplectic CM centralizer. Its ring is
$\mathcal O=\mathbb Z_3[\zeta_3]$, and its group is
\[
 \{z\in1+3\mathcal O:N_{\mathbb Q_3(\zeta_3)/\mathbb Q_3}(z)=1\}.
\]
The convergent logarithm identifies this group with
$3(\zeta_3-\zeta_3^2)\mathbb Z_3$.
The element
\[
          \epsilon=\frac{1+3\zeta_3}{1+3\zeta_3^2}           \tag{B.40.20}
\]
is a topological generator, since
\[
 \frac{\log\epsilon}{3(\zeta_3-\zeta_3^2)}\equiv1\pmod3.
\]
Indeed the first logarithm term gives one; terms whose
index is divisible by three vanish, and all remaining
terms of index at least two are divisible by three.

The full glued action of $\epsilon$ is exactly that
of the actual arithmetic element $\tau\overline\tau^{-1}$
from B.37. Its cyclotomic character is one. Its actual
outer arithmetic lift commutes with $\mathcal A$ over
$K$, so it fixes $\xi$.
Continuity then shows that the entire strict symplectic
centralizer fixes $\xi$.
Consequently any intertwiner satisfying (B.40.18)
gives the same transported geometric extension class.
This proves that the computation does not select an
unjustified Betti-to-local marking.

The comparison with the characteristic-two curve is
functorial for the cover (B.40.7) and all its displayed
automorphisms. The cover has degree nine, prime to two,
and extends with the smooth marked curve.
The prime-to-two specialization theorem
\citep[XIII, \S\S2.9--2.12]{BSGA1} therefore applies.
The tangent changes in this comparison are inner paths
already killed in (B.40.3).

\par\smallskip\noindent\textbf{Two congruences select the actual candidate.}\enspace 
We now apply (B.40.4) to all the candidates already
proved exhaustive in B.36:
\[
                 e=b+s h_3+t h_9,\qquad
                    s\in\mathbb Z/3, t\in\mathbb Z/9.
                                                               \tag{B.40.21}
\]
Write $e_4(s,t)\in Q/DQ$ for their projected restrictions
to $\mathbb F_4$, formed using the actual map $1+F_2$
as in B.36. Thus the thirteen-coordinate vector $e$ in
(B.40.21) and its rank-$320$ representative $e_4$ are
kept distinct.
There is an exact small-denominator formula for the
rational inverse needed here. The four eigenvalues
of $D$ give
\[
 D^4-6D^3-144D^2+54D+1215I=0,\qquad
 81D^{-1}=-\frac{D^3-6D^2-144D+54I}{15}.                       \tag{B.40.22}
\]
The numerator applied to each $e_4$ is divisible by three;
the remaining denominator five is a unit.
The resulting $81D^{-1}e_4$ is integral, consistently with
the largest local Smith exponent four.

Let $K_{\rm inv}$ be the four-column integral invariant
lattice after (B.40.19). Select its unit minor in
quotient rows $(46,51,80,82)$.
For
\[
                 y=81D^{-1}e_4-81h\pmod {81}
\]
condition (B.40.5) is precisely membership of $y$ in
the image of $K_{\rm inv}$ modulo $81$.
Eliminating the four invariant coordinates gives two
particularly short necessary equations. The first row
has coefficients
\[
 (51:-6, 54:1, 80:-27, 82:39),
\]
and the second has coefficients
\[
 (46:-15, 51:-17, 80:-11, 81:1, 82:33).
\]
All entries are taken modulo $81$; omitted coordinates
are zero. In the adapted Hall basis these quotient
positions are, respectively,
\[
\begin{array}{c|rrrrrr}
 \text{position}&46&51&54&80&81&82\\ \hline
 \text{Hall word}&(4,0,0)&(4,0,5)&(4,0,8)&(5,0,0)&(5,0,1)&(5,0,2).
\end{array}
\]
Both rows annihilate the invariant lattice modulo $81$.
Their evaluations on $y$ for (B.40.21) are
\[
                  27s,\qquad 27+9t\pmod {81}.                \tag{B.40.23}
\]
Thus $s=0$ and $t=6$ are forced.
Substitution makes every one of the $320$ membership
equations zero, proving sufficiency as well.
This yields exactly (B.40.1), and reading its twelfth
coordinate gives (B.40.2).

The complete executed calculation is in
\texttt{b\_\allowbreak elliptic\_\allowbreak schreier.py},
\texttt{b\_\allowbreak geometric\_\allowbreak invariants.py}, and
\texttt{b\_\allowbreak geometric\_\allowbreak phase.py}; their JSON and compressed
matrix files retain the words, basis changes, all
geometric displacements, invariant lattices, and both
full sparse rows.
The independent verifier
\texttt{c\_\allowbreak geometric\_\allowbreak phase\_\allowbreak verify.py}
recomputes all forty reversed words, uses a different
strict symplectic intertwiner with seed $(1,3)$,
uses the exact integer invariant lattice, and solves
$D$ by DVR elimination rather than (B.40.22).
It obtains the same unique thirteen-coordinate vector.
The working modulus is $3^{10}$, and the largest
DVR pivot valuation in the rational inverse is four.
Thus at least six digits remain after that division;
the compact gluing loses only one further digit.
In particular every displayed congruence modulo $81$
is strictly below the retained precision. All word
operations, integral basis changes, and the independent
invariant-lattice average are exact. The conclusion
is an exact finite-module equality.

Finally, restriction from $\mathbb Q_2$ to its unramified
quadratic extension is injective on three-primary
cohomology, since corestriction composed with restriction
is multiplication by two. Thus the selected restriction
determines the asserted class over $\mathbb Q_2$.
B.38 proves that the global formal refinements differ
by exactly its displayed $C_3$, and (B.40.6) proves
that its two nonzero differences change the actual
geometric class. Hence precisely one global refinement
has the actual geometric extension.
Existence is supplied by the actual fundamental-group
action, and (B.40.1) explicitly selects its phase.
This completes the integral refinement of the
nonzero rational Ross class already identified in
B.20 and B.37.
\end{proof}
